\clearpage
\setcounter{aufgabe}{0}
\einheitskopf{Einheit 1 von 5 · Prozente als Anteile}\label{e1}

\begin{aufgabe}{Prozente als Anteile – am Streifen. Der ganze Streifen sind 100\,\%. Wie viel Prozent sind grau? Beispiel: Sind 3 der 10 Kästchen grau, so sind das 30\,\%.}
\streifenreihe{a/20,b/60,c/10,d/90}
\begin{teilezwei}
\tz \leerfeld[\%] & \tz \leerfeld[\%] \\
\tz \leerfeld[\%] & \tz \leerfeld[\%] \\
\end{teilezwei}
Färbe im Streifen grau ein.
\begin{teile}
\teil 40\,\% des Streifens
\teil 65\,\% des Streifens
\end{teile}
\streifenreihe{e/0,f/0}
\end{aufgabe}

\begin{aufgabe}{Prozente als Anteile – umrechnen. Schreibe als Prozentangabe.}
\beispiel{\frac{17}{100} &= 17\,\%}
\begin{teilezwei}
\tz $\frac{23}{100} = \leerfeld[\%]$ & \tz $\frac{9}{100} = \leerfeld[\%]$ \\
\tz $\frac{57}{100} = \leerfeld[\%]$ & \tz $\frac{80}{100} = \leerfeld[\%]$ \\
\tz $\frac{13}{25} = \leerfeld[\%]$ & \tz $0{,}36 = \leerfeld[\%]$ \\
\end{teilezwei}
Schreibe als Dezimalzahl.
\begin{teile}
\teil $62\,\% = \leerfeld$
\end{teile}
Schreibe als Bruch und kürze so weit wie möglich.
\begin{teile}
\teil $40\,\% = \leerfeld$
\end{teile}
Schreibe als Prozentangabe.
\begin{teile}
\teil In einer Lieferung ist jedes zehnte Glas zerbrochen. \leerfeld[\%]
\end{teile}
Auf einer Packung steht: 16\,\% der Schrauben sind zu kurz. Kreuze die Aussage an, die dasselbe bedeutet.
\begin{teile}
\teil \kreuz{Jede 16.\ Schraube ist zu kurz.}\\
\kreuz{16 von 100 Schrauben sind zu kurz.}\\
\kreuz{In der Packung sind 16 Schrauben zu kurz.} \hfill (P10 2020)
\end{teile}
\end{aufgabe}

\begin{aufgabe}{Prozente als Anteile – ordnen.}
\begin{teile}
\teil Ordne der Größe nach, mit der kleinsten Angabe zuerst: 8\,\%, 80\,\%, 18\,\%. \\ \leerfeld\leerfeld\leerfeld
\teil Ordne der Größe nach, mit der kleinsten Angabe zuerst: $0{,}3$, $25\,\%$, $\frac{1}{5}$. \\ \leerfeld\leerfeld\leerfeld
\teil In Packung A sind $\frac{2}{5}$ der Bonbons rot, in Packung B sind es $45\,\%$. In welcher Packung ist der Anteil roter Bonbons größer? \leerfeld
\end{teile}
\end{aufgabe}

\begin{aufgabe}{Anteile zu 100\,\% ergänzen. Wie viel Prozent fehlen noch zu 100\,\%?}
\beispiel{100\,\% - 35\,\% &= 65\,\%}
\begin{teilezwei}
\tz $40\,\% \rightarrow \leerfeld[\%]$ & \tz $25\,\% \rightarrow \leerfeld[\%]$ \\
\tz $70\,\% \rightarrow \leerfeld[\%]$ & \\
\end{teilezwei}
\begin{teile}
\teil Von den Gästen einer Feier sind $30\,\%$ Kinder und $45\,\%$ Frauen. Wie viel Prozent sind Männer? \leerfeld[\%]
\teil Die Tabelle zeigt, wie sich der Umsatz eines Kiosks auf die Getränke verteilt. Ergänze den fehlenden Anteil und gib an, welches Getränk den größten Anteil hat. \hfill (P10 2021)
\end{teile}
\sachtabelle{lc}{Getränk & Anteil am Umsatz}{Wasser & 28\,\%\\ Saft & 22\,\%\\ Limonade & \leerzelle\\ Tee & 15\,\%}
\begin{teile}
\teil Größter Anteil: \leerfeld
\end{teile}
\end{aufgabe}

\begin{aufgabe}{Prozente als Anteile – Fehler finden. Auf einer Kiste steht: „Jeder achte Becher ist beschädigt.“ Jonas schreibt dafür 8\,\%.}
\begin{teile}
\teil Benenne den Fehler und gib den richtigen Prozentsatz an. \\ Fehler: \dsleer\dsleer\dsleer \quad richtig: \leerfeld[\%]
\teil In einer zweiten Kiste ist jeder zwanzigste Becher beschädigt. Wie viel Prozent sind das? \leerfeld[\%]
\end{teile}
\end{aufgabe}

\begin{aufgabe}{Prozente als Anteile – Begründen. Begründe ohne Taschenrechner.}
\begin{teile}
\teil Warum ist $\frac{1}{20}$ dasselbe wie $5\,\%$? \\ \dsleer\dsleer\dsleer\dsleer
\teil Lea sagt: „$0{,}5\,\%$ ist dasselbe wie $0{,}5$.“ Hat sie recht? Begründe. \\ \dsleer\dsleer\dsleer\dsleer
\end{teile}
\end{aufgabe}

\begin{aufgabe}{Anteilsaussage prüfen und korrigieren. Die Tabelle zeigt, wie die Schüler einer Schule zum Unterricht kommen.}
\sachtabelle{lcccc}{Weg & zu Fuß & Fahrrad & Bus & Auto}{Anteil & 35\,\% & 25\,\% & 30\,\% & 10\,\%}
In der Schülerzeitung stehen zwei Sätze dazu. Einer davon ist falsch.
\begin{teile}
\teil Kreuze den falschen Satz an. \\
\kreuz{Ein Viertel der Schüler kommt mit dem Fahrrad.}\\
\kreuz{Ein Drittel der Schüler kommt mit dem Bus.}
\teil Schreibe den falschen Satz richtig auf. \hfill (P10 2023) \\ \dsleer\dsleer\dsleer\dsleer\dsleer
\end{teile}
\end{aufgabe}
